\song{Pískající cikán (Irská / Spirituál kvintet)}
%Mara_2017_01_01

\ve{}
\ch{G}Dívka \ch{Ami}loudá se \ch{Hmi}vini - \ch{Ami}cí, \ \ \ch{G}tam, kde \ch{Ami}zídka je \ch{Hmi}níz - \ch{Ami}ká,\\
\ch{G}tam, kde \ch{Ami}stráň končí \ch{Hmi}voní\ch{C}cí, \ch{G}si písnič\ch{C}ku někdo \ch{G}pí\ch{C}sk\ch{D}{á}.

\cho{}
\ch{G Ami Hmi Ami G Ami Hmi Ami}{\textit{(pískání)}}\\
\ch{G Ami Hmi C G C G C D}{\textit{(pískání)}}

\ve{}
Ohlédne se a \uv{Propána!}, v stínu, kde stojí líska,\\
švárného vidí cikána, jak leží, písničku píská.

\cho{}

\ve{}
Chvíli tam stojí potichu, písnička si jí získá,\\
domů jdou spolu do smíchu, je slyšet cikán, jak píská.

\cho{}

\ve{}
Jenže tatík, jak vidí cikána, pěstí do stolu tříská,\\
\uv{Ať táhne pryč, vesta odraná, groš nemá, něco ti spíská.}

\cho{}

\ve{}
Teď smutnou dceru má u vrátek, jen Bůh ví, jak se jí stýská,\\
\uv{kéž vrátí se mi zas nazpátek ten který v dálce si píská.}

\ve{}Pár šídel honí se po louce, v trávě rosa se blýská,\\
cikán, rozmarýn v klobouce, jde dál a písničku píská.

\cho{}

\ve{}
Na závěr zbývá už jenom říct, v čem je ten kousek štístka:\\
peníze často nejsou nic, má víc, kdo po svém si píská.

\esong